\documentclass[11pt,a4paper]{article}
\usepackage[utf8]{inputenc}
\usepackage{amsmath,amssymb,amsthm}
\usepackage{listings}
\usepackage{geometry}
\usepackage{hyperref}
\geometry{margin=2.5cm}

\newtheorem{theorem}{Theorem}
\newtheorem{lemma}[theorem]{Lemma}

\lstset{
  language=C++,
  basicstyle=\ttfamily\small,
  keywordstyle=\bfseries,
  breaklines=true,
  frame=single,
  numbers=left,
  numberstyle=\tiny,
  tabsize=4,
  showstringspaces=false
}

\title{IOI 2008 -- Islands}
\author{Solution}
\date{}

\begin{document}
\maketitle

\section{Problem Statement}
Given a \textbf{functional graph} (each node has exactly one outgoing edge with a positive weight), find the sum of the longest simple paths within each connected component.

\section{Solution}

\subsection{Structure of Functional Graphs}
Each connected component of a functional graph contains exactly one cycle. Nodes not on the cycle form trees rooted at cycle nodes (the ``rho'' shape).

\subsection{Algorithm}
For each connected component:
\begin{enumerate}
    \item \textbf{Find the cycle.} Walk from any unvisited node, following outgoing edges, until revisiting a node.
    \item \textbf{Compute tree depths.} For each cycle node $c$, compute the maximum depth $d_c$ in the tree hanging off $c$ (via DFS on reverse edges, excluding cycle edges). Also track the tree diameter (longest path within each tree).
    \item \textbf{Combine paths through the cycle.} The longest path either:
    \begin{itemize}
        \item lies entirely within one tree (the tree diameter), or
        \item passes through part of the cycle: $d_i + d_j + \text{dist}(c_i, c_j)$.
    \end{itemize}
\end{enumerate}

\subsection{Doubling Trick for Circular Optimization}
For $L$ cycle nodes with depths $d_0, \ldots, d_{L-1}$ and cumulative edge weights $\text{prefix}[k]$, we want:
\[
\max_{\substack{0 \le i < j < L \\ j - i < L}} \bigl(d_i + d_j + \text{prefix}[j] - \text{prefix}[i]\bigr).
\]
By doubling the cycle array to length $2L$ and using a sliding-window maximum (deque) with window size $L$, we compute:
\[
\max_j \Bigl(d_{j \bmod L} + \text{prefix}[j] + \max_{i \in [j-L+1,\; j-1]} \bigl(d_{i \bmod L} - \text{prefix}[i]\bigr)\Bigr)
\]
in $O(L)$ time.

\subsection{Complexity}
\begin{itemize}
    \item \textbf{Time:} $O(N)$ total.
    \item \textbf{Space:} $O(N)$.
\end{itemize}

\section{C++ Solution}

\begin{lstlisting}
#include <bits/stdc++.h>
using namespace std;

int main(){
    ios::sync_with_stdio(false);
    cin.tie(nullptr);

    int N;
    cin >> N;

    vector<int> nxt(N + 1);
    vector<long long> w(N + 1);
    for (int i = 1; i <= N; i++)
        cin >> nxt[i] >> w[i];

    // Reverse adjacency list for tree DFS
    vector<vector<pair<int, long long>>> radj(N + 1);
    for (int i = 1; i <= N; i++)
        radj[nxt[i]].push_back({i, w[i]});

    vector<bool> visited(N + 1, false);
    vector<bool> onCycle(N + 1, false);
    long long totalAns = 0;

    for (int start = 1; start <= N; start++) {
        if (visited[start]) continue;

        // Find the cycle in this component
        vector<int> path;
        unordered_set<int> inPath;
        int cur = start;
        while (!inPath.count(cur) && !visited[cur]) {
            inPath.insert(cur);
            path.push_back(cur);
            cur = nxt[cur];
        }

        if (visited[cur]) {
            for (int u : path) visited[u] = true;
            continue;
        }

        // Extract cycle starting at 'cur'
        vector<int> cycle;
        vector<long long> cycleW;
        bool found = false;
        for (int u : path) {
            if (u == cur) found = true;
            if (found) {
                cycle.push_back(u);
                onCycle[u] = true;
            }
        }
        int L = cycle.size();
        cycleW.resize(L);
        for (int i = 0; i < L; i++)
            cycleW[i] = w[cycle[i]]; // edge cycle[i] -> cycle[(i+1)%L]

        for (int u : path) visited[u] = true;

        // For each cycle node, DFS into its tree to find max depth and diameter
        vector<long long> depth(L, 0);
        long long compAns = 0;

        for (int ci = 0; ci < L; ci++) {
            int root = cycle[ci];
            long long treeDiam = 0;

            function<long long(int)> treeDFS = [&](int u) -> long long {
                long long mx = 0;
                for (auto [v, wt] : radj[u]) {
                    if (onCycle[v]) continue;
                    if (visited[v]) continue; // safety
                    visited[v] = true;
                    long long childDepth = treeDFS(v) + wt;
                    treeDiam = max(treeDiam, mx + childDepth);
                    mx = max(mx, childDepth);
                }
                return mx;
            };

            depth[ci] = treeDFS(root);
            compAns = max(compAns, treeDiam);
        }

        // Sliding window maximum on doubled cycle
        vector<long long> prefix(2 * L + 1, 0);
        for (int i = 0; i < 2 * L; i++)
            prefix[i + 1] = prefix[i] + cycleW[i % L];

        deque<int> dq;
        for (int j = 0; j < 2 * L; j++) {
            while (!dq.empty() && dq.front() <= j - L)
                dq.pop_front();
            if (!dq.empty()) {
                int i = dq.front();
                long long val = depth[j % L] + prefix[j]
                              + depth[i % L] - prefix[i];
                compAns = max(compAns, val);
            }
            long long jVal = depth[j % L] - prefix[j];
            while (!dq.empty()) {
                long long backVal = depth[dq.back() % L] - prefix[dq.back()];
                if (backVal <= jVal) dq.pop_back();
                else break;
            }
            dq.push_back(j);
        }

        totalAns += compAns;
    }

    cout << totalAns << "\n";
    return 0;
}
\end{lstlisting}

\section{Notes}
The functional-graph decomposition (cycle detection + tree DFS) is a fundamental technique. The sliding-window maximum on the doubled cycle array handles the circular distance computation in $O(L)$, and the constraint $j - i < L$ ensures the path does not revisit any cycle node.

A subtle point: the tree DFS uses the reverse adjacency list and must avoid crossing back onto the cycle. The \texttt{visited} flag and \texttt{onCycle} flag together enforce this.

\end{document}
